% Chapter 16: Why the World Keeps Oscillating
\chapter{Why the World Keeps Oscillating}

\step{A Counterintuitive Opening}

Consider three familiar little events you probably seldom question.

First, a swing. Someone beside it gives a tiny push each time it reaches its highest point, barely more than a touch. Minutes later it swings higher than anyone else's. Much stronger but poorly timed pushes, sometimes against its return, instead slow it quickly.

Second, the kitchen. Hang a stainless-steel spoon from a finger, let it swing like a pendulum, and time it with your phone. You may notice something odd: wide or narrow swings take almost the same time per round trip. Wider swings seem more hurried, narrower ones more leisurely, but the period refuses to follow intuition.

Third, a bridge. In 1940, Washington's Tacoma Narrows Bridge twisted like a rope in a wind that was not especially unusual, then broke and fell into the strait, all captured on film. How can a steel bridge vibrate like a guitar string?

One thread connects them. Objects displaced from where they settle often do not return quietly, but move back and forth for some time. A plucked guitar string sounds, a nudged pendulum swings, a car crossing a speed bump bounces several times, and water ripples when a cup is tapped. Even lattice atoms, microphone circuits, and atomic nuclei in MRI machines oscillate at their own rhythms. Why does the world love oscillation? We will answer, and explain how a gentle push can send a swing soaring.

Emphasize this first: you have encountered all these phenomena. We introduce no inaccessible new facts, only a shared framework. Next time a clothesline shakes in wind, a refrigerator compressor hums, or a subway handhold sways during acceleration, you will recognize another oscillator.

\step{Make a Guess First}

Commit your intuition to these questions before continuing, with no changing your bets.

\begin{enumerate}[leftmargin=2em]
  \item Replace a swing's rider with someone twice as heavy. Does the round-trip time become noticeably longer, shorter, or nearly unchanged?
  \item To make a swing rise higher, is push strength or timing more important?
  \item A mass hangs from a spring and oscillates vertically after being pulled down and released. Move it to the Moon, where gravity is one sixth of Earth's. Does its rhythm change?
  \item Why does a car bounce several times after a pothole instead of sinking once and stopping? Would removing its shock absorbers produce more bouncing or less?
\end{enumerate}

Write your answers and grade them at the end. Most people miss at least one, often the chapter's central question. A preview: none needs a new law, only Newton's laws and a new viewpoint. Physics often advances by organizing old phenomena differently rather than finding new ones. Oscillation is one such organizing idea.

\step{Hands-on Experiment}

\begin{experiment}[A Ruler at the Table Edge and Keys on a Spring]
Experiment 1: hold a steel or plastic ruler firmly against a table with part projecting over the edge. Bend the free end downward slightly and release it. The ruler vibrates and buzzes. Change the overhanging length: shorter means faster vibration and higher pitch; longer means slower and lower. Without any instrument, you discover that vibration has its own rhythm, determined by the system, the ruler and exposed length, not how hard you bend it. A stronger bend gives more vigorous vibration but nearly the same rhythm.

Experiment 2: hang keys from a rubber band or thin spring. Hold it still, then pull the keys gently down and release. They bounce for a long time. Count oscillations in ten seconds. Remove half the keys and count again. Some intuitions predict lighter means slower, others faster. Let your own measurements decide.

Experiment 3: let the keys hang still, then move your hand up and down by a tiny amount. Rapid random shaking does little. Slow down toward the natural rhythm you just measured. The instant the rhythms match, the keys suddenly leap up and down even though your hand moves less than before. Remember: matching rhythm matters much more than strength. For another step, deliberately slow your hand slightly after matching. The large motion gradually fades; the energy savings account opens only at particular rhythms.
\end{experiment}

These formula-free experiments already contain every keyword: equilibrium position, natural rhythm, amplitude, and matching rhythm. We will give them formal names.

If possible, replay slow-motion video of the keys and notice two things. Near the lowest point they seem to rush past, while at the high and low endpoints they linger: speed is greatest in the middle and zero at the ends. Also, large and small amplitudes give nearly the same rhythm. These direct observations provide the formulas' entire empirical foundation. A good theory explains what you saw; it does not invent facts.

\step{The Old Model Runs into Trouble}

We know Newton's laws: net force changes momentum; zero net force maintains rest or uniform straight-line motion. They explain each instant of the keys' motion. Pull them down and spring tension exceeds weight, giving upward acceleration. At the middle net force vanishes. Beyond it, net force points down, slowing, stopping, and reversing them.

Yet following instant after instant feels like chasing a car we can never catch. We want a panoramic map: what is the whole system doing? Why does it not stop where net force vanishes and all forces balance?

Intuition says balanced forces should make an object settle. Experiment says the opposite: it passes equilibrium fastest, at its least settled. Why not stop? Inertia: velocity does not vanish automatically when force disappears. After overshooting, the spring compresses or stretches again and net force reverses, slowing and returning it. Inertia prevents stopping in time; restoring force prevents escaping far. Together they make back-and-forth motion inevitable.

This is the first conceptual shift: oscillation needs no strange new law. It is the tug-of-war between inertia and a force pulling toward equilibrium. Where both exist, oscillation appears. Almost every stable equilibrium develops a restoring force under small displacement: a ball in a bowl, a taut string, compressed gas, a bent ruler, even molecular chemical bonds. Hence an oscillating world. Distinguish the roles: disturbed systems often oscillate is an observed fact; restoring force proportional to displacement is a mathematical approximation near equilibrium; inertia battling restoration is our physical explanation. None implies an object wants to go home.

Ask the reverse: what does not oscillate? A bowl of still water has a flat surface; tap it with a chopstick and it ripples before settling. Still water and jelly differ not in whether they move, but how they return after a push. Jelly's elasticity restores its shape and produces vibration. Water returns through gravity and surface tension and also ripples. Honey is so viscous that energy dissipates quickly and it creeps back without oscillating. Whether motion oscillates, and for how long, depends on stiffness, inertia, and dissipation. This viewpoint predicts an unfamiliar system's general behavior more readily than chasing forces instant by instant.

\step{Inventing New Concepts}

Turn our everyday expressions into precise concepts, as physicists do.

\textbf{Equilibrium position}: where an undisturbed system rests. For spring-hung keys it is their resting height; for a swing, the bottom; for a ruler, its unbent shape. What matters is that displacement produces a force pulling back. Equilibrium need not mean nothing is happening: the hanging spring is already stretched enough for elasticity to balance gravity. We care not about absolute extension but distance from equilibrium. Move the coordinate origin there and gravity's constant contribution is settled in advance, leaving restoration to keep the remaining account. This small change of origin makes the formulas clean.

\textbf{Restoring force}: the net force always directed toward equilibrium. It is not a new kind of force beside gravity and elasticity, but a role net force plays. Elasticity and weight jointly supply it for hanging springs; gravity's component along the swinging path supplies it for a pendulum. Like centripetal force, the name describes the job, not the source.

\textbf{Simple harmonic motion}: if restoring force is exactly proportional to displacement from equilibrium,
\[
  F = -kx,
\]
the resulting oscillation is simple harmonic. Here \(x\) is displacement and \(k\) the proportionality constant; a spring's \(k\) is its familiar spring constant. The minus sign is the formula's soul: force always opposes displacement, left when you move right and right when you move left. Springs obey this relationship, Hooke's law, very accurately when neither stretched nor compressed too far.

\textbf{Amplitude, period, frequency, and phase}: greatest distance from equilibrium is amplitude \(A\), how wide the swing is. Round-trip time is period \(T\), how long one cycle takes. Cycles per unit time give frequency \(f=1/T\), measured in hertz, Hz, cycles per second. Phase says where you are within the cycle. Two swings of identical amplitude and frequency may lead or lag one another: different phases. Abstract though phase sounds, you use it daily. Pushing a swing at the right moment means matching its phase.

These terms make everyday devices readable. A tuning fork marked 440 Hz completes 440 cycles per second, the standard musical pitch. An accurately cut quartz crystal in a watch vibrates electrically at 32768 Hz, two to the fifteenth power. Halving the rate fifteen times gives one pulse per second, the watch's tick. A phone's vibration motor spins a small off-center weight, providing hundreds of forced vibrations per second. MRI makes hydrogen nuclei in the body wobble at a particular frequency in a magnetic field, then pushes their swings with radio waves at that frequency. Watches and bodies share this language because both have equilibria and restoring forces.

\begin{figure}[htbp]
\centering
\begin{tikzpicture}[thick,>=Latex]
  % Wall
  \draw[pattern=north east lines] (0,0) rectangle (0.25,2.4);
  \draw (0.25,0) -- (0.25,2.4);
  % Spring, shown as a coil
  \draw[decorate,decoration={coil,aspect=0.5,segment length=4pt,amplitude=4pt}] (0.25,1.2) -- (3.3,1.2);
  % Block
  \draw[fill=blue!15] (3.3,0.8) rectangle (4.3,1.6);
  \node at (3.8,1.2) {$m$};
  % Ground
  \draw (2.8,0.6) -- (5.6,0.6);
  % Equilibrium line
  \draw[dashed,gray] (3.8,0.2) -- (3.8,2.4);
  \node[gray,below] at (3.8,0.2) {Equilibrium};
  % Displacement arrow
  \draw[<->,red] (3.8,2.15) -- (5.0,2.15) node[midway,above]{$x$};
  % Restoring-force arrow
  \draw[->,red,very thick] (4.3,1.2) -- (3.3,1.2) node[midway,below]{$F=-kx$};
\end{tikzpicture}
\caption{A spring oscillator displaced by \(x\) experiences a net force toward equilibrium; greater displacement means greater restoring force.}
\end{figure}

\step{The Model in One Sentence}

\begin{oneline}
When displacement produces a net force pulling back toward equilibrium roughly in proportion to that displacement, inertia and restoration make the system oscillate at a fixed rhythm; that rhythm depends only on its stiffness, restoring coefficient \(k\), and inertia, mass \(m\), not on how hard you push.
\end{oneline}

Check the pieces. Inertia explains overshooting equilibrium; restoration explains returning; a system-defined rhythm explains why stronger ruler bending changes loudness rather than pitch. Stubbornness and laziness are metaphors for spring stiffness and mass, but unlike vague adjectives they have precise meanings: \(k\) has units and can be measured.

Also state what the model is not. It is not a belief that everything has a spirit and wants to return home. Restoring force is structural mechanics: stretch a spring, displace atoms from their lowest-energy spacing, and force appears. Nor does oscillation normally have to stop. The ideal model oscillates forever; stopping requires energy leakage, an extra explanation rather than a default. Automatic stopping is an Aristotelian intuition we dismantled in kinematics and must dismantle again here.

\step{The Mathematical Translator}

Intuition says a stiffer spring pulls harder and oscillates faster, while a heavier object responds sluggishly and oscillates slower. Translate these proportionalities into the harmonic rhythm:
\[
  \omega=\sqrt{\frac{k}{m}},\qquad
  T=2\pi\sqrt{\frac{m}{k}},\qquad
  f=\frac{1}{2\pi}\sqrt{\frac{k}{m}}.
\]
The angular frequency \(\omega\), in radians per second, differs from ordinary frequency by \(2\pi\): \(\omega=2\pi f\). Why square roots rather than direct proportionality? Doubling \(k\) doubles force and acceleration, making the same round trip quicker. But acceleration acts through two time accumulations: acceleration builds velocity, and velocity builds displacement. Their combination reduces time only to \(1/\sqrt{2}\) of its original value. The mathematical bridge gives the rigorous origin.

Test the new formulas with special cases:

\begin{itemize}[leftmargin=2em]
  \item \(k\to 0\), a spring so soft it effectively vanishes, gives \(T\to\infty\). The object drifts away without returning. No restoring force, no oscillation: sensible.
  \item \(m\to\infty\), an infinitely massive load, gives \(T\to\infty\). Nothing can set infinite inertia oscillating: sensible.
  \item Amplitude \(A\) is absent. Double amplitude and period stays unchanged. The spoon's surprising behavior is now a prediction. This holds only when the restoring force is strictly proportional to displacement, a central limitation later.
\end{itemize}

A pendulum, swing, or hanging spoon is approximately harmonic at small angles. Before its formula, guess with dimensions. Period is measured in seconds. The relevant pendulum parameters are length \(L\) in meters, mass \(m\) in kilograms, and gravitational acceleration \(g\) in meters per second squared. To construct seconds, kilograms have nowhere to cancel, so mass must be absent. Dividing \(L\) by \(g\) gives seconds squared; take the square root. The almost unique possibility is \(T\sim\sqrt{L/g}\). This is not proof, only dimensional consistency, but already predicts mass independence and square-root length dependence. Guessing with units and then testing or deriving is one of physicists' favorite shortcuts.

The rigorous derivation in the mathematical bridge gives
\[
  T=2\pi\sqrt{\frac{L}{g}},
\]
with length \(L\) and gravitational acceleration \(g\). Check \(g\to 0\), as in a space station: \(T\to\infty\), and the pendulum stops functioning, a point for our thought experiment. Quadruple \(L\) and the period only doubles. Doubling swing length therefore slows the rhythm by about forty percent, not a factor of two. Mass \(m\) never appears: our first prediction answer is nearly unchanged period for double body weight.

\textbf{An estimate with real numbers.} For length \(L=1.0\) meters and Earth's \(g\approx9.8\) meters per second squared,
\[
  T=2\pi\sqrt{\frac{1.0}{9.8}}\approx 2.0\ \text{s}.
\]
A round trip takes about two seconds. This is the famous seconds pendulum of clockmaking: roughly one meter long, ticking once each half-period. Traditional standing clocks use this length. Conversely, measure length with a ruler and period with a phone to measure local gravitational acceleration at home within a few percent. Oscillation links time and force, which is why it makes a clock.

\textbf{Another real-data estimate: car suspension.} A typical car weighs about 1,500 kilograms, or 375 per wheel. When a person gets in, its body sinks about two centimeters. From \(k=mg/x\), estimate each suspension spring's stiffness as \(375\times9.8/0.02\approx1.8\times10^{5}\) newtons per meter. Then
\[
  T=2\pi\sqrt{\frac{375}{1.8\times10^{5}}}\approx 0.3\ \text{s},
\]
giving a body-bounce natural frequency around three hertz, though actual designs are lower, around one or two. This explains the one-or-two-bounces-per-second rhythm after a speed bump, set by car and springs rather than driving speed. It also explains sensitivity to low-frequency motion in carsickness: abdominal organs are themselves oscillators with nearby natural frequencies. Matching external drive makes the organs resonate, naturally causing discomfort.

Now energy. Release an oscillator at its amplitude: speed is zero and energy is entirely elastic or gravitational potential, \(\frac12 kA^{2}\). At equilibrium the spring returns to its original shape, potential vanishes, and all energy becomes kinetic, \(\frac12 mv^{2}\), with maximum speed. At the far endpoint it becomes potential again. Oscillation repeatedly exchanges stored and moving energy while the total \(\frac12 kA^{2}\) stays fixed. Double amplitude and energy quadruples, so raising a swing a little requires more energy than it appears.

Energy also explains the slow-motion observation. With a fixed total, less potential means more kinetic: equilibrium has lowest potential and highest speed; endpoints hold all the potential and force speed to zero. It explains why pushing at the bottom is most efficient: speed is greatest, so the same force over the same distance transfers the most energy. Near the top the swing nearly stops and a long push stores little. Energy accounting often saves more work than a force diagram; we will repeatedly use it.

\begin{mathbridge}
Newton's second law for a spring oscillator gives
\[
  m\frac{d^{2}x}{dt^{2}}=-kx
  \quad\Longrightarrow\quad
  \frac{d^{2}x}{dt^{2}}+\omega^{2}x=0,
  \qquad \omega^{2}=\frac{k}{m}.
\]
This second-order differential equation asks for a function whose second derivative is proportional to its negative. Sine and cosine qualify: let \(x(t)=A\cos(\omega t+\varphi)\), differentiate twice, and obtain \(x''=-\omega^{2}x\). The square root appears because angular frequency drops out twice during differentiation. Amplitude \(A\) and initial phase \(\varphi\) come not from the equation but from initial conditions: how far you pull and whether release includes initial velocity. Mathematics thus explains system-defined rhythm and user-defined amplitude. The exact pendulum equation, \(\theta''=-(g/L)\sin\theta\), is not harmonic. Only at small angles, replacing \(\sin\theta\approx\theta\), does it become so. That is the small-angle approximation's origin. For energy, \(\frac12 mv^{2}+\frac12 kx^{2}=\frac12 kA^{2}=\text{constant}\) holds at every instant for an ideal undamped oscillator.
\end{mathbridge}

\begin{figure}[htbp]
\centering
\begin{tikzpicture}[thick,>=Latex,scale=1.0]
  \draw[->] (0,0) -- (7.2,0) node[right]{$t$};
  \draw[->] (0,-1.6) -- (0,1.9) node[above]{$x$};
  \draw[blue,thick,domain=0:6.9,samples=100] plot (\x,{1.2*cos(\x*180)});
  \draw[dashed,gray] (0,1.2) -- (6.9,1.2);
  \draw[dashed,gray] (0,-1.2) -- (6.9,-1.2);
  \draw[<->,red] (6.45,0) -- (6.45,1.2) node[midway,right]{$A$};
  \draw[<->,red] (1.0,-1.5) -- (3.0,-1.5) node[midway,below]{$T$};
  \fill (1.0,1.2) circle (1.2pt); \fill (3.0,1.2) circle (1.2pt);
\end{tikzpicture}
\caption{Harmonic displacement versus time is a sine or cosine curve: neighboring peaks are one period \(T\) apart, and peak distance from equilibrium is amplitude \(A\).}
\end{figure}

\step{The Model's Limits}

The harmonic model is clean and elegant, but three explicit conditions limit it.

\textbf{Condition 1: restoring force must be proportional to displacement.} Overstretch a spring and permanent deformation invalidates Hooke's law. A pendulum's \(\sin\theta\approx\theta\) works only at small angles. Beyond thirty or forty degrees, period lengthens noticeably; at \(60^{\circ}\), it is about seven percent longer than the small-angle result. Amplitude-independent period becomes an approximation. Intuition easily errs here: small everyday angles give small errors, so we promote an approximation to a law.

\textbf{Condition 2: no energy leakage.} Real friction and air resistance turn energy into heat and reduce amplitude over time: damping. The keys stop not because their strength runs out, but because mechanical energy becomes internal energy. Damping is not always bad. Car shock absorbers deliberately drain spring oscillations, or a pothole would leave you bouncing until dark. Correct a common misconception: springs alone do not make a comfortable car. Springs cushion the impact; dampers finish the job. A spring-only car truly bounces.

Distinguish three damping levels. Too little gives many cycles before stopping, like a plucked guitar string. At a critical value, the oscillator returns without overshooting: critical damping. Good door closers and ammeter needles aim for this fast, nonoscillatory return. Too much damping resembles falling into honey, creeping slowly back without even one oscillation. The same spring changes character with damping. Shock-absorber tuning adjusts exactly this balance.

\textbf{Condition 3: resonance is not a universal explanation.} Textbooks often call Tacoma Narrows a bridge destroyed by wind resonance, but closer analysis finds torsional motion more like aerodynamic self-excitation, flutter. Bridge twist alters airflow, which amplifies the twist, even in constant wind. It shares the swing's core idea, continuous energy input at the right phase, but is subtler than ordinary resonance. Blaming everything on resonance misuses the concept. The real engineering lesson is that sustained energy input can destabilize an oscillating system.

\begin{figure}[htbp]
\centering
\begin{tikzpicture}[thick,>=Latex,scale=1.0]
  \draw[->] (0,0) -- (7.0,0) node[right]{Drive frequency};
  \draw[->] (0,0) -- (0,3.1) node[above]{Amplitude};
  \draw[blue,thick] plot[smooth] coordinates {(0,0.5) (1.5,0.7) (2.6,1.4) (3.2,2.8) (3.8,1.4) (5.0,0.8) (6.5,0.55)};
  \draw[red,thick] plot[smooth] coordinates {(0,0.45) (1.5,0.6) (2.6,1.0) (3.2,1.5) (3.8,1.0) (5.0,0.7) (6.5,0.5)};
  \draw[dashed,gray] (3.2,0) -- (3.2,2.8);
  \node[gray,below] at (3.2,0) {$f_0$};
  \node[blue] at (4.6,2.4) {Low damping: sharp peak};
  \node[red] at (5.6,1.2) {High damping: broad peak};
\end{tikzpicture}
\caption{Forced-oscillation resonance: amplitude peaks near the natural frequency \(f_0\); less damping makes the peak taller and sharper.}
\end{figure}

\begin{challenge}
Forced oscillation obeys \(x''+2\beta x'+\omega_0^2 x=(F_0/m)\cos\omega t\), where \(\beta\) describes damping and \(\omega\) is drive frequency. Its steady state oscillates at that frequency with amplitude
\[
  A(\omega)=\frac{F_0/m}{\sqrt{(\omega_0^{2}-\omega^{2})^{2}+4\beta^{2}\omega^{2}}}.
\]
Check two limits: as \(\omega\to 0\), \(A\to F_0/k\), so a slow push becomes static stretching. As \(\beta\to 0\) and \(\omega\to\omega_0\), the denominator vanishes and amplitude diverges. Undamped resonance is mathematically unbounded; reality means something eventually breaks first. Engineers characterize peak sharpness by quality factor \(Q=\omega_0/(2\beta)\). A tuning fork's \(Q\) can exceed a thousand, ringing long after one strike. Car suspension deliberately has \(Q\) near one, giving only one bounce per pothole. The same mathematics governs RLC circuits and laser cavities: one oscillation structure, many physical carriers.
\end{challenge}

\step{Explaining the Opening Phenomenon}

Now grade the opening's three puzzles.

\textbf{Why gentle pushes raise a swing.} A swing approximately behaves as a pendulum with natural period \(T=2\pi\sqrt{L/g}\), nearly amplitude-independent at small angles. Gentle pushes are a small drive. Match the natural rhythm and every push aligns with motion, depositing net energy each cycle and gradually increasing amplitude: resonance. Weak force is fine if each cycle deposits into the same account. Random pushes sometimes oppose the return, withdrawing energy and preventing growth. This also answers prediction 2: timing, not strength, is decisive.

\textbf{Why the spoon keeps its rhythm.} At small angles, \(T=2\pi\sqrt{L/g}\) contains no amplitude. Wider swings travel farther, but gravity's tangential component grows proportionally. The effects cancel and period stays constant. This is isochronism, which led Galileo, timing with his pulse, to use pendulums for clocks. Intuition equated farther travel with longer time and forgot the growing force.

\textbf{Why a bridge vibrates apart.} A bridge deck is not perfectly rigid and has natural frequencies. Steady wind sheds alternating vortices around it, supplying continuous drive. Near a natural frequency, resonance amplifies motion; the aerodynamic self-excitation described earlier adds to it until structural limits are exceeded. Subsequent suspension-bridge designs require wind-tunnel tests, keeping natural frequencies away from common wind-driving frequencies and using damping to absorb energy. Music shows the benefit of the same principle: guitar bodies and violin soundboards respond near string frequencies, amplifying tiny string vibrations into audible sound. Resonance makes music or destroys bridges according to where we want energy stored.

Skyscrapers face the same problem and related solutions. Near Taipei 101's eighty-eighth floor hangs a steel ball about 5.5 meters across and 660 tonnes in mass, suspended by cables like a giant pendulum. When typhoon winds drive the building near its natural frequency, engineers tune the ball to that rhythm and make it oppose the building. As the tower sways east, the ball lags, and connecting hydraulic dampers continually turn its vibration energy into heat. This tuned mass damper recruits destructive resonance as a bodyguard: since energy must go somewhere, arrange a safe destination. Many supertall buildings in earthquake-prone Japan and windy coastal regions conceal similar stabilizing masses.

\step{Thought Experiments and Challenges}

\begin{thought}[Weighing Yourself in a Space Station]
In weightlessness, a bathroom scale reads zero, not because you lack mass but because its spring-compression-until-support-balances-weight principle fails. How do astronauts monitor body mass? Seat them in a chair connected to a known spring, pull it from equilibrium, release, and measure period \(T\). Then \(m=kT^{2}/(4\pi^{2})\) directly gives mass. Space stations really use such devices. The deeper point: this measures inertia, not gravitational weight. Inertial mass remains fully meaningful in weightlessness. A pendulum fails there, \(g\to0\) and \(T\to\infty\), while a spring oscillator works normally. Their dependence on gravity differs completely. Would the spring chair give a different result in distant interstellar space?
\end{thought}

Oscillation also bridges to later chapters. A guitar string is countless neighboring oscillators pulling one another; their collective dance is a wave, Chapter 17's subject. Complicated vibrations split into simple-frequency components, the entrance to Chapter 18's Fourier ideas. Farther ahead, chemical-bond springs hold lattice atoms near equilibrium, and their thermal vibrations relate closely to temperature. Quantum harmonic oscillators have energy in discrete portions and can never be completely still; zero-point energy awaits in the quantum section. Circuits qualify too: a capacitor--inductor loop exchanges electric and magnetic energy like kinetic and potential energy, a genuine oscillator that lets radios select stations. Understand this one spring and you hold a key to much of physics.

\begin{puzzles}
\item Move the same spring and mass from Earth to the Moon, where gravity is about one sixth as strong. How does its vertical oscillation period change? What if you move a pendulum of unchanged length instead? (Hint: check which period formula contains \(g\). Spring stiffness \(k\) belongs to the spring itself.)
\item Skilled swing riders crouch at the bottom and stand at the top, rising higher without anyone pushing. Where do they steal the energy? (Hint: standing raises the center of mass and does work against the downward force; crouching at maximum speed costs little. The effective length changes each cycle, parametric driving rather than ordinary force driving.)
\item An old radio's tuning knob adjusts the natural frequency of an electrical capacitor--inductor oscillator. Why does one station become audible while others fall quiet, and why do nearby stations sometimes interfere? (Hint: inspect resonance curves. Overlap means an overly broad peak, too much damping and poor selectivity.)
\item A poorly designed bridge sways increasingly when walkers' step frequency approaches its natural frequency. London's Millennium Bridge really closed for repairs two days after opening because of sway. The remedy added dampers rather than thicker steel beams. Why can making the bridge dissipate more energy work better than making it stronger? (Hint: resonant amplitude balances input and dissipated power. Reduce \(Q\) and the peak collapses.)
\end{puzzles}
